\documentclass[10pt]{article}
\usepackage[landscape,margin=0.5in]{geometry}
\usepackage{graphicx}
\usepackage{caption}
\usepackage{parskip}
\pagestyle{empty}
\captionsetup{font=footnotesize, labelformat=empty, skip=2pt}

\newcommand{\panel}[2]{%
  \begin{minipage}[t]{0.235\textwidth}\centering
    \includegraphics[width=\linewidth,height=3.4cm,keepaspectratio]{figures/#1}\\[1pt]
    {\footnotesize #2}
  \end{minipage}\hfill}

\begin{document}
\begin{center}
{\Large\textbf{Visual Changelog --- HAnim LOA5 Demonstration}}\\[1pt]
{\small How the running figure and classroom skeleton evolved. Each panel is a
verified render; commit ids let you check out that exact state in git.}
\end{center}
\vspace{4pt}

{\large\textbf{1.\ The runner}}\\[3pt]
\panel{fig_stick_loa4.png}{\textbf{a) Generated stick figure (LOA-4).}
Real HAnim joints, our gait, gait synced to ground speed.\\ \texttt{b1a4541}}
\panel{fig_canonical_skeleton.png}{\textbf{b) Canonical LOA5 skeleton.}
The Web3D \texttt{AllBonesLOA5} figure: 150 joints, 243 bone meshes.\\ \texttt{47ffc4f}}
\panel{fig_runner_bonemesh.png}{\textbf{c) Bone-mesh runner.}
The real skeleton runs the slalom course, driven through standard joints.\\ \texttt{47ffc4f}}
\panel{fig_runner_xite_course.png}{\textbf{d) Browser delivery (X\_ITE).}
The real HAnim scene in a web player---no flattening.\\ \texttt{0259f0c}}

\vspace{10pt}
{\large\textbf{2.\ The classroom}}\\[3pt]
\panel{fig_flux_blackboard.png}{\textbf{e) Generated textures.}
Chalkboard \& posters made locally with FLUX.1, applied as \texttt{ImageTexture}.\\ \texttt{132cb9d}}
\panel{fig_pbr_torso.png}{\textbf{f) PBR look-dev.}
Bone meshes re-authored to \texttt{PhysicalMaterial}; form returns under IBL.\\ \texttt{a05e930}}
\panel{fig_classroom_buttons.png}{\textbf{g) Interactive classroom.}
LOA5 skeleton on a stand; chalkboard buttons switch Walk/Run/Jump.\\ \texttt{b26a291}}
\panel{fig_buttons_closeup.png}{\textbf{h) Chalk-styled controls.}
Buttons restyled to read as chalk drawings on the board.\\ \texttt{9d0c601}}

\vspace{14pt}
\begin{center}\footnotesize
Browse any step: \texttt{git show <id> -{}-stat} (what changed) \quad$\cdot$\quad
\texttt{git checkout <id>} then \texttt{git checkout main} (time-travel and back).
\end{center}
\end{document}
